% Chapter 6 -- Claim: High-Interest Currencies Depreciate Enough to Offset Their Yield
% Requires: amsmath, amssymb, booktabs, array, tabularx, longtable, graphicx, tcolorbox (breakable, skins), hyperref or url.
% Defines the callout box only if the preamble has not.
\providecolor{fenavy}{HTML}{1B2A6B}
\providecolor{feshade}{HTML}{F3F5FA}
\providecolor{feorange}{HTML}{D9822B}
\makeatletter
\@ifundefined{callout}{%
  \newtcolorbox{callout}{breakable, enhanced,
    colback=white, colframe=gray!70, boxrule=0.5pt, arc=1mm,
    borderline west={2.5pt}{0pt}{gray!70},
    left=8pt, right=8pt, top=4pt, bottom=4pt}%
}{}
\makeatother

\providecolor{fenavy}{HTML}{1B2A6B}
\providecolor{feshade}{HTML}{F3F5FA}
\makeatletter
% Shaded box for optional technical material that general readers may skip.
\@ifundefined{technote}{%
  \newtcolorbox{technote}{breakable, enhanced,
    colback=feshade, colframe=fenavy!60, boxrule=0.5pt, arc=1mm,
    left=8pt, right=8pt, top=4pt, bottom=4pt}%
}{}
\makeatother

\chapter{Claim: High-Interest Currencies Depreciate Enough to Offset Their Yield}
\label{ch:11}
\providecommand{\term}[1]{\emph{#1}}

Banks in some countries pay far more interest on savings than banks in others. A saver might reasonably ask why she should not simply keep her money wherever interest is highest. The catch is that the money must be converted into another currency and back again, and exchange rates move. This chapter asks whether the exchange rate quietly takes back the extra interest, as economic theory says it should.

When the Bretton Woods system of fixed exchange rates, the postwar arrangement in which major currencies were pegged to the U.S. dollar, ended in 1973, currencies began to float, and economists gained a new laboratory. Suppose investors care only about expected returns. Then a currency that pays a higher interest rate must be expected to fall by just enough to cancel its extra yield. Otherwise everyone would borrow in the low-rate currency and lend in the high-rate one until the opportunity disappeared. This idea is called \term{uncovered interest parity}\index{uncovered interest parity (UIP)|textbf}, and much of international macroeconomics rests on it. Fama (1984) tested it with the floating exchange rates of the 1970s. Instead of finding that high-interest currencies fell, as the theory predicts, he found that they tended to rise.

This result became known as the \term{forward premium puzzle}\index{forward premium puzzle}, one of the most persistent findings in finance. It has a practical counterpart in the \term{carry trade}\index{carry trade}: borrowing in a low-interest currency to invest in a high-interest one. In 2007 a Japanese household could borrow yen at less than 1\% and deposit Australian dollars at more than 6\%. Japanese retail investors, nicknamed ``Mrs.\ Watanabe,'' did so on a large scale, alongside hedge funds. Later research asked why the trade pays. Brunnermeier, Nagel and Pedersen (2008) showed that it suffers sudden crashes when the borrowed currency, such as the yen, jumps in value. To see whether high-interest currencies share a common risk, Lustig, Roussanov and Verdelhan (2011) grouped currencies by their interest rates and found a common risk that explains much of the return.



% USD/JPY spot and implied one-year forward rate, 2016--2026.
% Values digitized from the source chart; replace with the original data when available.
% Requires: pgfplots, pgfplotstable, textcomp (for \textyen).
\providecolor{fenavy}{HTML}{1B2A6B}
\providecolor{feorange}{HTML}{D9822B}
\providecolor{fered}{HTML}{B2182B}
\pgfplotstableread{
  year spot fwd
  2016.890 113.88 112.44
  2017.220 110.80 109.51
  2017.484 111.80 110.16
  2017.968 112.09 109.61
  2018.232 105.73 103.35
  2018.485 110.21 107.13
  2018.738 112.99 109.71
  2018.991 109.11 106.18
  2019.233 110.21 107.13
  2019.475 107.22 105.24
  2019.992 108.02 106.13
  2020.223 107.03 106.58
  2020.344 107.03 106.83
  2020.454 107.42 106.83
  2020.619 106.23 105.83
  2020.762 104.84 104.54
  2020.894 103.65 103.30
  2020.960 103.25 102.65
  2021.048 103.95 104.49
  2021.180 108.81 107.72
  2021.279 110.40 109.91
  2021.334 110.21 109.91
  2021.411 110.60 110.21
  2021.719 110.80 110.50
  2021.807 111.70 111.05
  2021.950 113.88 113.78
  2022.038 115.77 114.78
  2022.210 121.70 119.35
  2022.500 135.94 132.11
  2022.731 144.29 138.27
  2022.790 150.15 143.00
  2023.006 130.87 125.71
  2023.226 132.46 126.50
  2023.490 144.09 136.73
  2023.732 149.06 141.70
  2023.820 151.34 143.80
  2023.974 142.40 135.99
  2024.238 151.24 144.54
  2024.315 158.89 150.74
  2024.480 160.98 154.10
  2024.656 147.37 142.10
  2024.986 156.90 151.40
  2025.261 149.65 144.48
  2025.492 144.19 139.72
  2025.745 147.86 143.39
  2025.999 156.71 152.73
  2026.252 156.61 152.53
  2026.384 157.90 153.63
  2026.450 159.29 155.56
  2026.575 163.86 159.10
  2026.681 158.69 155.22
  2026.747 156.81 153.23
  2026.890 159.80 155.28
}\usdjpydata

\begin{figure}[htbp]
\centering
\scalebox{0.65}{%
\begin{tikzpicture}[
  note/.style={draw=gray!70, fill=white, rounded corners=2pt, font=\small,
               align=left, inner sep=3pt},
  pointer/.style={-stealth, thin},
]
\begin{axis}[
  scale only axis, width=1.02\linewidth, height=0.5\linewidth,
  xmin=2016.6, xmax=2027.0, ymin=90, ymax=172,
  xtick={2017,2018,...,2026},
  xticklabel style={/pgf/number format/1000 sep={}},
  ytick={100,110,...,170},
  ylabel={Japanese yen per U.S. dollar (\textyen)},
  grid=major, major grid style={gray!25}, set layers,
  legend style={at={(0.5,-0.1)}, anchor=north, font=\small, draw=gray!60,
                legend columns=2, cells={anchor=west},
                /tikz/every even column/.append style={column sep=1.2em}},
  clip=false,
]
  % Regime shading
  \begin{pgfonlayer}{axis background}
    \fill[fenavy!7]   (axis cs:2016.6,90) rectangle (axis cs:2022.0,172);
    \fill[feorange!9] (axis cs:2022.0,90) rectangle (axis cs:2026.95,172);
  \end{pgfonlayer}
  \addlegendimage{area legend, fill=fenavy!7, draw=fenavy!25}
  \addlegendentry{2016--2021: stable range (\textyen103--\textyen115)}
  \addlegendimage{area legend, fill=feorange!9, draw=feorange!30}
  \addlegendentry{2022--2026: weaker yen, wide forward discount}

  % Series
  \addplot[fenavy, very thick] table[x=year, y=spot] {\usdjpydata};
  \addlegendentry{Spot USD/JPY (\textyen\ per \$1)}
  \addplot[feorange, very thick, dashed] table[x=year, y=fwd] {\usdjpydata};
  \addlegendentry{Implied 1-year forward (covered interest parity)}

  % Annotations
  \node[note, anchor=north east] (n1) at (axis cs:2020.5,101.5) {\textbf{Dec 2020 low}\\\textyen103.25};
  \draw[pointer] (n1.east) -- (axis cs:2020.93,103.0);
  \node[note, anchor=east] (n2) at (axis cs:2021.55,126) {\textbf{Mar 2022}\\\textbf{Fed hikes begin}\\(\textyen121.70)};
  \draw[pointer] (n2.east) -- (axis cs:2022.17,121.4);
  \node[note, anchor=east] (n3) at (axis cs:2022.25,152.5) {\textbf{Oct 2022}\\\textbf{BoJ intervention}\\(\textyen150.15)};
  \draw[pointer] (n3.east) -- (axis cs:2022.75,150.3);
  \node[note, draw=fered, text=fered, anchor=east] (n4) at (axis cs:2026.1,167) {\textbf{Jul 29, 2026 peak}\\\textyen163.86};
  \draw[pointer, fered] (n4.east) -- (axis cs:2026.53,164.0);

  % End labels
  \node[fenavy, anchor=west, font=\small\bfseries] at (axis cs:2027.05,160.6) {Spot: \textyen159.80};
  \node[feorange, anchor=west, font=\small\bfseries] at (axis cs:2027.05,154.6) {1Y Fwd: \textyen155.28};
\end{axis}
\end{tikzpicture}%
}
\caption{When U.S. interest rates are higher than Japan's, the dollar's forward rate sits below its spot rate. USD/JPY spot rate and implied one-year forward rate, 2016--2026. 
It quotes the exchange rate as yen per dollar, the usual market convention, which is the inverse of the $S$ used in this chapter. A rise in the plotted line therefore means that the dollar has appreciated and the yen has depreciated, so $S$ has fallen.
The forward rate comes from covered interest parity. It lies below the spot rate because U.S. interest rates exceed Japanese rates, and the gap widened after the Federal Reserve began raising rates in 2022.}
\label{fig:06-usdjpy}
\end{figure}


The answer matters to central banks, companies and investors. If interest parity fails, a central bank that raises rates may see its currency rise rather than fall. A company that leaves its foreign revenue unhedged is taking a predictable bet. And an investor in the carry trade is being paid to bear crash risk, not collecting a free lunch. The claim on trial\index{claims-on-trial method} is the textbook one: \textbf{high-interest currencies fall enough to cancel their extra yield, so the carry trade earns nothing on average.}

The lab in Section~\ref{sec:06-lab} repeats Fama's (1984) test for nine currencies against the dollar, using free exchange rates and interest rates from FRED. It then builds the carry portfolios of Lustig, Roussanov and Verdelhan (2011). The goal is to check whether uncovered interest parity holds today, and what the carry trade's returns and crashes look like.

By the end of the chapter you will be able to:

\begin{itemize}
  \item compute a forward exchange rate from covered interest parity\index{covered interest parity (CIP)}, and the profit or loss on a carry trade;
  \item run and interpret the Fama regression, and explain the forward premium puzzle\index{forward premium puzzle};
  \item compare the explanations for the puzzle and the evidence for each;
  \item judge what a currency carry fund is paid for and how it loses money.
\end{itemize}


\section{Foundations: exchange rates, forwards and interest parity}
\label{sec:06-foundations}

\subsection{Exchange rates}

The \term{spot rate}\index{spot rate|textbf} $S$ is the price of one currency in terms of another. In this chapter $S$ is the number of U.S. dollars per unit of foreign currency, so a rise in $S$ means the foreign currency has \term{appreciated}, or become more valuable. Changes are measured in logs: $\Delta s_{t+1} = \ln(S_{t+1}/S_t)$.
It’s called the spot rate because the trade is settled “on the spot,” meaning right away rather than on some agreed future date. The word comes from the old commercial phrase for a cash sale with immediate delivery.



\subsection{Forward rates}

A \term{forward rate}\index{forward rate|textbf} $F$ is an exchange rate agreed today for a currency trade that will take place on a specific future date. It is quoted in the same units as the spot rate $S$, and it serves as the market's break-even benchmark for future currency conversions.
“Spot” and “forward” are a pair: the spot rate is for exchanging currencies now, and the forward rate is for exchanging them later.

A forward contract has two sides. The party that agrees to deliver a currency on the future date \textbf{sells it forward}, and the party that agrees to receive it \textbf{buys it forward}. Every currency trade swaps one currency for another, so selling dollars forward for yen is the same as buying yen forward. A Japanese exporter that will be paid in dollars sells dollars forward. A U.S. importer that owes a Japanese supplier yen buys yen forward.

Companies use forward rates to remove the currency risk from payments they already know they will make or receive. Figure~\ref{fig:06-usdjpy} gives an example. It quotes the exchange rate as yen per dollar, the usual market convention, which is the inverse of the $S$ used in this chapter. At the end of the sample, the spot rate was \textyen159.80 per dollar and the one-year forward rate was \textyen155.28. Suppose a Japanese car exporter expects to receive \$10 million from U.S. customers in a year. By selling the dollars forward today, it locks in \textyen1,552.8 million, whatever happens to the exchange rate. If it does not hedge, its yen income depends on the spot rate a year later. In 2024 the dollar fell from about \textyen161 to about \textyen147 within a few months. A fall of that size would leave the exporter with only \textyen1,470 million, about \textyen83 million less than the hedge guarantees. If instead the dollar stays at \textyen159.80, the exporter would have received \textyen1,598 million without the hedge, so the hedge costs it \textyen45.2 million. \textbf{Hedging is therefore not free, but its cost is known in advance.} Why the forward rate sits below the spot rate is the subject of the next subsection.


\subsection{Covered interest parity (CIP)}

An investor with one dollar has two safe ways to invest it for a year. On the first route, she deposits the dollar at the dollar interest rate $i$ and ends the year with
\begin{equation}
  1 + i \text{ dollars.}
  \label{eq:06-cip-home}
\end{equation}
On the second route, she invests abroad and hedges the currency risk. Because $S$ is the number of dollars per unit of foreign currency, her dollar buys $1/S$ units of foreign currency today. She deposits them at the foreign interest rate $i^*$, so at the end of the year she holds $(1+i^*)/S$ units of foreign currency. She has also agreed today to convert them back at the forward rate $F$, so she ends the year with
\begin{equation}
  \frac{F(1+i^*)}{S} \text{ dollars.}
  \label{eq:06-cip-abroad}
\end{equation}

Every rate on both routes is known today, so neither route carries any currency risk. If one route paid more than the other, traders could borrow through the cheaper route and invest through the better one for a riskless profit. \term{Arbitrage}\index{arbitrage} therefore forces the two payoffs to be equal. Setting Equation~\eqref{eq:06-cip-home} equal to Equation~\eqref{eq:06-cip-abroad} and rearranging gives \term{covered interest parity}\index{covered interest parity|textbf}:
\begin{equation}
  \frac{F}{S} = \frac{1 + i}{1 + i^*}.
  \label{eq:06-cip}
\end{equation}

In words, covered interest parity says that the forward rate differs from today's rate by exactly the gap between the two countries' interest rates, so hedging a foreign deposit removes its extra interest. The shaded box shows the algebra.

\begin{technote}
\textbf{Technical note: covered interest parity in logs.} The condition is easier to read in logs. Let $f = \ln F$ and $s = \ln S$. Taking logs of both sides of Equation~\eqref{eq:06-cip} and using the approximation $\ln(1+x) \approx x$, which is accurate when $x$ is small, gives
\begin{equation}
  f - s = \ln(1+i) - \ln(1+i^*) \approx i - i^*.
  \label{eq:06-cip-log}
\end{equation}
The left side is approximately the percentage gap between the forward rate and the spot rate, and the right side is the gap between the two interest rates. Covered interest parity therefore says that the forward rate differs from the spot rate by exactly the interest-rate differential. When the dollar rate is higher than the foreign rate, $F$ is above $S$, so the dollar buys less foreign currency in the forward market than in the spot market. This shortfall, called the \term{forward discount}\index{forward discount|textbf} on the dollar, takes back the extra interest that dollar deposits earn. The interest rates $i$ and $i^*$ must cover the same period as the forward contract, so a one-year forward rate is compared with one-year interest rates.
\end{technote}


Figure~\ref{fig:06-usdjpy} shows the parity at work. The ratio $159.80/155.28 \approx 1.029$ implies that the one-year dollar rate exceeded the yen rate by about 2.9 percentage points. An investor who holds dollars earns that extra interest, and the lower forward rate takes it back when the dollars are converted, so hedged dollars and hedged yen pay the same. The figure also shows how the gap follows interest rates. From 2016 to 2021, when interest rates in both countries were close to zero, the forward rate stayed close to the spot rate. After the Federal Reserve began raising rates in March 2022, while Japanese rates stayed far lower, the gap widened to several yen.

Covered interest parity held closely until 2008, and it has shown small but persistent deviations since then (Du, Tepper and Verdelhan 2018).

If the investor does not hedge, she must convert her foreign currency back at the unknown spot rate a year from now, and \textbf{uncovered interest parity}\index{uncovered interest parity|textbf} is the hypothesis that this unhedged foreign investment earns, on average, the same return as the dollar deposit.

\subsection{Uncovered interest parity (UIP)}

\textbf{Uncovered interest parity (UIP)}\index{uncovered interest parity (UIP)|textbf} goes one step further. It says that the spot rate \emph{expected} in the future equals today's forward rate, so a high-interest currency is expected to fall by its interest advantage:
\begin{equation}
  E_t[\Delta s_{t+1}] = i_t - i^*_t
  \label{eq:06-uip}
\end{equation}
Unlike covered interest parity, UIP is not enforced by arbitrage. Holding a foreign currency without a forward contract is risky, so UIP holds only if investors demand no extra return for that risk.

Figure~\ref{fig:06-uip} shows the difference with the numbers of the worked example below. An investor with \$1 million can earn the U.S. rate of 1\% and end the year with \$1,010,000 for certain. Or she can buy Australian dollars, earn 5\%, and convert back at whatever the spot rate turns out to be. With a forward contract, the second route would also pay \$1,010,000 for certain; that is covered interest parity. Without one, its dollar payoff depends on the future spot rate $S_1$, which nobody knows today. UIP says that investors expect $S_1$ to equal the forward rate, \$0.673, so that the two routes have the same \emph{expected} payoff. Only the first route, however, is free of risk.

\begin{figure}[htbp]
\centering
\scalebox{0.8}{%
\begin{tikzpicture}[
  box/.style={draw=fenavy, fill=fenavy!6, rounded corners=3pt, line width=0.8pt,
              align=center, inner sep=6pt, text width=4.6cm},
  plain/.style={draw=black, fill=white, rounded corners=3pt, line width=0.8pt,
              align=center, inner sep=6pt, text width=4.6cm},
  risky/.style={draw=feorange, fill=feorange!8, rounded corners=3pt, line width=0.8pt,
              dashed, align=center, inner sep=6pt, text width=4.6cm},
  wide/.style={draw=fenavy, fill=fenavy!6, rounded corners=3pt, line width=0.8pt,
              align=center, inner sep=7pt, text width=12cm},
  arr/.style={-stealth, line width=0.9pt},
]
  % Start
  \node[box, text width=5.2cm] (start) at (0,0) {\large\bfseries \$1,000,000 today};

  % Left route: U.S. deposit
  \node[plain] (usd) at (-3.5,-1.8) {\textbf{U.S. deposit}\\Earn 1\% for 1 year};
  \node[plain] (usdend) at (-3.5,-7.95)
    {\textbf{\$1,010,000}\\known today, no risk};

  % Right route: unhedged Australian deposit
  \node[box] (buy) at (3.5,-1.8) {\textbf{Buy A\$ at spot, \$0.70}\\A\$1,428,571.43};
  \node[box] (earn) at (3.5,-3.55) {\textbf{Earn 5\%}\\A\$1,500,000};
  \node[risky] (conv) at (3.5,-5.45)
    {\textbf{No forward contract}\\Sell A\$ at next year's\\spot rate $S_1$, unknown today};
  \node[risky] (audend) at (3.5,-7.95)
    {\textbf{A\$1,500,000} $\times\, S_1$\\[2pt]
     \small unchanged at 0.70: \$1,050,000\\
     \small down 10\% to 0.63: \$945,000};

  % Bottom: the UIP condition
  \node[wide] (uip) at (0,-10.1)
    {\textbf{Uncovered interest parity:} both routes have the same \emph{expected} payoff\\[4pt]
     $E[S_1] = 0.70 \times \dfrac{1.01}{1.05} \approx 0.673 = F$
     \quad$\Rightarrow$\quad expected payoff $\approx$ \textbf{\$1,010,000}};

  % Arrows
  \draw[arr] (start.south) -- ++(0,-0.35) -| (usd.north);
  \draw[arr, fenavy] (start.south) -- ++(0,-0.35) -| (buy.north);
  \draw[arr] (usd.south) -- node[fill=white, inner sep=3pt] {Hold for 1 year} (usdend.north);
  \draw[arr, fenavy] (buy) -- (earn);
  \draw[arr, fenavy] (earn) -- (conv);
  \draw[arr, feorange] (conv) -- (audend);
  \draw[arr, gray] (usdend.south) -- (usdend.south |- uip.north);
  \draw[arr, gray] (audend.south) -- (audend.south |- uip.north);
\end{tikzpicture}%
}
\caption{Uncovered interest parity says the risky foreign route and the safe home route should pay the same on average. Uncovered interest parity, with the numbers of the worked example: one year, a spot rate of \$0.70 per Australian dollar, 1\% in U.S. dollars and 5\% in Australian dollars. The Australian route is not hedged, so its dollar payoff depends on next year's spot rate $S_1$. UIP says that investors expect the Australian dollar to fall to the forward rate of about \$0.673, so the two routes have the same expected payoff. Only the U.S. route is certain.}
\label{fig:06-uip}
\end{figure}


Can an investor profit by borrowing in a low-interest currency and investing in a high-interest one, or will the high-interest currency depreciate by just enough to wipe out the gain, as uncovered interest parity predicts? This strategy is called the \textbf{carry trade}\index{carry trade|textbf}.
\subsection{The carry trade}

Can an investor profit by borrowing in a low-interest currency and investing in a high-interest one, or will the high-interest currency depreciate by just enough to wipe out the gain, as uncovered interest parity predicts? This strategy is called the \textbf{carry trade}\index{carry trade|textbf}. The investor converts the borrowed money into the high-interest currency and invests it there. At the end of the period, she converts the proceeds back without any hedge and repays the loan. Measured in dollars, her excess return\index{return!excess} is the change in the exchange rate minus the interest-rate differential:
\begin{equation}
  rx_{t+1} = \Delta s_{t+1} - (i_t - i^*_t)
  \label{eq:06-carry}
\end{equation}
Here $s$ is the log of the number of dollars per unit of foreign currency, so $\Delta s_{t+1} > 0$ means that the foreign currency has appreciated. The interest rates $i_t$ and $i^*_t$ are known when the trade is placed at time $t$. When the foreign rate is the higher one, $i_t - i^*_t$ is negative, so the trade earns the interest gap and then gains or loses whatever the exchange rate does. Under UIP its expected value is zero, because the currency is expected to fall by exactly the interest gap. Put another way, the carry trade buys high-interest currencies at a discount in the forward market and hopes they will not fall by as much as the discount.

\textbf{Worked example.} The U.S. one-year rate is 1\%, the Australian rate is 5\%, and the spot rate is \$0.70 per Australian dollar (Figure~\ref{fig:06-carry}).

\begin{itemize}
  \item Covered interest parity sets the one-year forward rate at $0.70 \times 1.01/1.05 \approx \$0.673$.
  \item Under UIP, the spot rate expected in a year is also \$0.673, a fall of 3.8\%.
  \item The carry trade borrows \$1 million at 1\%, buys A\$1.429 million, and invests it at 5\% to receive A\$1.5 million in a year. Its result depends on the spot rate then:
  \begin{itemize}
    \item if the rate is unchanged at \$0.70, the proceeds are \$1.05 million, the loan costs \$1.01 million, and the profit is \$40,000;
    \item if the rate falls to \$0.673, as UIP predicts, the proceeds are \$1.01 million and the profit is zero;
    \item if the rate falls 10\% to \$0.63, the proceeds are \$945,000 and the loss is \$65,000.
  \end{itemize}
\end{itemize}


\begin{figure}[htbp]
\centering
\scalebox{0.8}{%
\begin{tikzpicture}[
  % Profit ($ thousands) = 1500 x (spot - 0.6733): +40k if unchanged at 0.70,
  % zero at the UIP forecast 0.673, -65k after a 10% fall to 0.63
  declare function={carry(\s) = 1500*(\s - 0.67333);}
]
\begin{axis}[
  width=0.61\linewidth, height=0.34\linewidth,
  axis lines=left, axis line style={-stealth},
  xmin=0.6, xmax=0.75, ymin=-110, ymax=120,
  xtick={0.6,0.62,...,0.74}, ytick={-100,0,100},
  xticklabel style={/pgf/number format/fixed, /pgf/number format/precision=2},
  xlabel={Spot rate in one year (\$ per Australian dollar)},
  ylabel={Carry profit (\$ thousands)},
  clip=false,
]
  \addplot[gray, thin, domain=0.6:0.75, samples=2] {0};
  \addplot[fenavy, very thick, domain=0.6:0.75, samples=2] {carry(x)};
  \fill[fenavy] (axis cs:0.70,40)  circle (2.6pt);
  \fill[fenavy] (axis cs:0.67333,0) circle (2.6pt);
  \fill[fenavy] (axis cs:0.63,-65) circle (2.6pt);
  \node[anchor=south east] at (axis cs:0.697,42) {unchanged: $+\$40$k};
  \node[anchor=north west] at (axis cs:0.676,-3) {UIP forecast 0.673: zero};
  \node[anchor=west] at (axis cs:0.633,-70) {$-10\%$: $-\$65$k};
\end{axis}
\end{tikzpicture}%
}
\caption{The carry trade profits if the high-interest currency holds its value and loses if it falls far enough. Profit on the worked example's carry trade against the spot rate a year later. The trade earns \$40,000 if the Australian dollar is unchanged, breaks even at the UIP forecast of \$0.673, and loses \$65,000 if the currency falls 10\%.}
\label{fig:06-carry}
\end{figure}

\section{Testing uncovered interest parity}
\label{sec:06-how}

UIP can be tested in two ways: one currency at a time, or across many currencies at once.


\subsection{The Fama regression.} 

The test is a regression of how much each currency actually moved on how much its interest rate differed from the dollar's. If UIP holds, the slope of that regression should be 1. The shaded box explains why.

\begin{technote}
\textbf{Technical note: from expectations to a regression.} Testing UIP runs into a measurement problem. UIP is a claim about the depreciation that investors expect:
\begin{equation}
  E_t[\Delta s_{t+1}] = i_t - i^*_t,
  \label{eq:06-uip-expected}
\end{equation}
but nobody can observe what the market expects. What can be observed is the actual change in the exchange rate a period later. One assumption bridges the gap. Under rational expectations, investors' forecasts can be wrong, but not systematically, so
\begin{equation}
  \Delta s_{t+1} = E_t[\Delta s_{t+1}] + e_{t+1},
  \label{eq:06-rational}
\end{equation}
where the forecast error $e_{t+1}$ cannot be predicted from anything known at time $t$, including the interest-rate differential. If UIP holds, substituting Equation~\eqref{eq:06-uip-expected} into Equation~\eqref{eq:06-rational} gives $\Delta s_{t+1} = (i_t - i^*_t) + e_{t+1}$. A regression of the actual change in the exchange rate on the differential,
\begin{equation}
  \Delta s_{t+1} = a + b\,(i_t - i^*_t) + e_{t+1},
  \label{eq:06-fama}
\end{equation}
should therefore find $a = 0$ and $b = 1$. The regression turns an untestable statement about expectations into a testable claim about data.
\end{technote}

Fama (1984) also had an important question in mind. Because $i_t - i^*_t \approx f_t - s_t$ by covered interest parity, the regression asks whether the forward rate is an unbiased forecast of the future spot rate, which was a central question about market efficiency. If the answer is no, the forward discount must contain something besides expected depreciation, namely a risk premium. Fama used the estimated $b$ to show that this premium must be large, must vary over time, and must move opposite to expected depreciation.

By covered interest parity, $i_t - i^*_t \approx f_t - s_t$, so this regression is the same as regressing the change in the exchange rate on the forward discount, which is how Fama (1984) wrote it. UIP implies $b = 1$: a foreign currency with a one-point interest advantage should fall by one point. If $b = 0$, exchange rates move at random, and the carry trade keeps the whole interest gap. If $b < 0$, high-interest currencies tend to \emph{rise}, and the carry trade earns even more than the interest gap (Figure~\ref{fig:06-fama}). This last case is the \textbf{forward premium puzzle}.

Substituting Equation~\eqref{eq:06-fama} into Equation~\eqref{eq:06-carry} links the regression directly to the carry trade's excess return:
\begin{equation}
  rx_{t+1} = a + (b - 1)\,(i_t - i^*_t) + e_{t+1}
  \label{eq:06-carry-fama}
\end{equation}
Under UIP, $a = 0$ and $b = 1$, so the interest-rate differential cannot predict the excess return. When $b < 1$, a foreign interest rate above the dollar rate, which makes $i_t - i^*_t$ negative, predicts a positive excess return, and the smaller $b$ is, the larger that return becomes.

\begin{figure}[htbp]
\centering
\scalebox{0.8}{%
\begin{tikzpicture}
\begin{axis}[
  width=0.46\linewidth, height=0.34\linewidth,
  axis lines=left, axis line style={-stealth},
  xmin=-0.5, xmax=4.5, ymin=-4.5, ymax=4.5,
  xtick={0,1,...,4}, ytick={-4,-2,0,2,4},
  xlabel={Foreign interest advantage, $i^* - i$ (\%)},
  ylabel={Expected change in foreign currency (\%)},
  clip=false,
]
  \addplot[gray, thin, domain=-0.5:4.5, samples=2] {0};
  % Typical estimate b = -0.9: the currency rises by 0.9 x its advantage
  \addplot[thick, densely dotted, domain=0:4.3, samples=2] {0.9*x};
  % Random walk b = 0
  \addplot[fegray, thick, dashed, domain=0:4.3, samples=2] {0};
  % UIP b = 1: falls by its interest advantage
  \addplot[fenavy, very thick, domain=0:4.3, samples=2] {-x};

  % Regions: what up and down mean
  \node[gray, anchor=west, font=\footnotesize] at (axis cs:-0.4,4.2)
    {foreign currency rises};
  \node[gray, anchor=west, font=\footnotesize] at (axis cs:-0.4,-4.2)
    {foreign currency falls};

  % Worked example at a 3-point interest advantage.
  % Carry return = expected change + interest advantage
  %              = vertical distance above the UIP line.
  \draw[gray, densely dotted] (axis cs:3,-4.5) -- (axis cs:3,4.5);
  \draw[-stealth, thick, gray] (axis cs:3,-3) -- (axis cs:3,2.55);
  \fill[fenavy] (axis cs:3,-3)  circle (2.2pt);
  \fill[fegray] (axis cs:3,0)   circle (2.2pt);
  \fill         (axis cs:3,2.7) circle (2.2pt);
  \node[anchor=south east, font=\footnotesize, align=right]
    at (axis cs:2.95,2.9) {carry: $3 + 2.7 = 5.7\%$};
  \node[anchor=south west, font=\footnotesize, fegray]
    at (axis cs:3.05,0.1) {carry: $3 + 0 = 3\%$};
  \node[anchor=north east, font=\footnotesize, fenavy]
    at (axis cs:2.95,-3.1) {carry: $3 - 3 = 0$};

  % Line labels
  \node[anchor=west, align=left] at (axis cs:5.3,3.87)
    {typical estimate\\($b \approx -0.9$):\\rises instead};
  \node[fegray, anchor=west, align=left] at (axis cs:5.3,0)
    {random walk\\($b = 0$)};
  \node[fenavy, anchor=west, align=left] at (axis cs:5.3,-4.3)
    {UIP ($b = 1$):\\falls by its\\interest advantage};
\end{axis}
\end{tikzpicture}%
}
\caption{In practice, high-interest currencies have tended to rise rather than fall, the opposite of what the theory predicts. Three readings of the Fama regression. The horizontal axis is $i^* - i$, the negative of the differential in Equation~\eqref{eq:06-fama}, so each line has slope $-b$. Under UIP ($b = 1$), a high-interest currency falls by its interest advantage. Under a random walk ($b = 0$), it does not move, and under the typical estimate ($b \approx -0.9$, Froot and Thaler 1990), it rises. The carry trade's expected return is the vertical distance above the UIP line. At an advantage of 3 points, it is zero under UIP, 3\% under a random walk, and 5.7\% under the typical estimate.}
\label{fig:06-fama}
\end{figure}

\textbf{Carry portfolios.} Each month, rank currencies by their interest rates, buy the highest-yielding ones and sell the lowest-yielding ones. If UIP holds, this portfolio earns nothing on average. A portfolio averages out the noise of individual currencies, and it is what investors actually trade.

\subsection{Why UIP might fail.} There are three broad explanations.

\begin{itemize}
  \item \textbf{A risk premium.} High-interest currencies tend to fall in global downturns, when investors flee to the dollar, the yen and the Swiss franc. Investors therefore demand extra return for holding them. Lustig, Roussanov and Verdelhan (2011) argue that carry returns are payment for exposure to this global risk.
  \item \textbf{Crash risk.} Carry returns may be steady most of the time and then collapse when carry traders all close their positions at once, as funding dries up. On this view the premium pays for those crashes (Brunnermeier, Nagel and Pedersen 2008).
  \item \textbf{Slow adjustment.} Investors may react slowly to changes in interest rates, so exchange rates move toward UIP only gradually. That would be a mispricing rather than a risk premium, and it should shrink as more money chases it.
\end{itemize}

The three explanations leave different fingerprints. A risk premium should show up as losses in bad times for the world economy. Crash risk should show up as negative skewness, with losses concentrated in a few episodes. Slow adjustment should show up as a premium that shrinks over time and does not depend on market conditions.



\section{The claim}
\label{sec:06-claim}

\begin{callout}
\textbf{Claim (uncovered interest parity).} The currency of a high-interest country is expected to fall by its interest advantage, so borrowing in low-interest currencies and lending in high-interest ones earns nothing on average.
\end{callout}

To test it, three points need to be settled:

\begin{itemize}
  \item \textbf{Over what period.} Over a few months, random exchange-rate moves swamp everything, so the test needs decades of data and many currencies.
  \item \textbf{Which currencies.} Nine developed-market currencies against the dollar. Emerging-market currencies offer larger interest gaps but also capital controls, default risk and much higher costs.
  \item \textbf{What would count against it.} Fama slopes significantly below 1 and a significant average carry return. The failure is deeper if the return survives costs and the slopes are below zero, not merely below one.
\end{itemize}

\section{The original evidence}
\label{sec:06-evidence}

\textbf{Fama (1984).} For nine currencies against the dollar after 1973, the slopes were not just below 1 but generally negative: high-interest currencies tended to rise. Fama showed that this implies a time-varying risk premium in forward rates. Across 75 published estimates, the average slope was about $-0.88$ (Froot and Thaler 1990).

\textbf{Carry as a risk factor.} Lustig, Roussanov and Verdelhan (2011) found that a portfolio long high-rate and short low-rate currencies\index{factor!HML (value)} earned substantial returns, and that a carry factor and a dollar factor explained much of the difference in returns across currencies. Menkhoff, Sarno, Schmeling and Schrimpf (2012) found that carry losses coincided with rises in global currency volatility\index{volatility}.

\textbf{Crashes.} Brunnermeier, Nagel and Pedersen (2008) showed that carry returns are negatively skewed\index{skewness}, with losses when risk appetite and funding\index{liquidity} fell, measured among other things by rises in the VIX\index{VIX}, as carry traders closed their positions together.

\textbf{Peso problems.}\index{peso problem|textbf} The premium might pay for rare disasters not yet in the data, a \term{peso problem}. Burnside, Eichenbaum, Kleshchelski and Rebelo (2011) protected carry positions with currency options and still found a substantial return, so unseen disasters could explain the premium only if priced far above their likely cost.

\section{Lab~\thechapter: Replicating Fama (1984) and Lustig, Roussanov and Verdelhan (2011)}
\label{sec:06-lab}

\textbf{Claim on trial\index{claims-on-trial method}:} high-interest currencies fall enough to cancel their extra yield. The lab replicates the Fama (1984) regression and the carry portfolios of Lustig, Roussanov and Verdelhan (2011), using free monthly data from FRED\index{FRED}: the Federal Reserve's exchange rates for nine developed-market currencies and OECD three-month interbank rates for each country and the United States. The interest gap stands in for the forward discount, which covered interest parity makes nearly equal, and FRED's monthly averages smooth returns and make crashes look milder.

\textbf{Part A, Fama regressions.} For each currency and for all together, the monthly change in the log exchange rate regressed on last month's U.S.-minus-foreign interest gap, with Newey--West\index{standard error!Newey--West (HAC)} standard errors (Lab Figure~\ref{ch:11}.1).

\textbf{Part B, the carry portfolio.} Each month, long the three highest-rate currencies and short the three lowest, split into interest earned and exchange-rate change (Lab Figure~\ref{ch:11}.2), with its distribution compared with a normal\index{distribution!normal} one (Lab Figure~\ref{ch:11}.3).

\textbf{Part C, robustness checks.} Three subperiods, costs\index{transaction costs} of 0 to 20 basis points\index{basis point}, exposure to stocks and to changes in the VIX, and the ten worst months.

\textbf{Part D, the investor's view.} Drawdowns\index{drawdown}, recovery times and carry returns against changes in the VIX (Lab Figures~\ref{ch:11}.4 and~\ref{ch:11}.5).

\label{sec:06-stress}
\textbf{Robustness checks.} Carry earns most of its return in long calm stretches and gives much back in a few months, so the start and end of the sample matter, and the period since 2008 matters most, because developed-country interest rates were close together for years. Major currencies have spreads\index{bid--ask spread} of a few basis points, so costs rarely eliminate the premium in developed markets. Traders call carry ``going up the stairs and down the elevator\index{skewness|textbf},'' and the lab's worst months should cluster in 1992, 1998, 2008 and August 2024. If carry is a risk premium\index{risk premium}, it should fall with stocks and with rises in the VIX. The lab ends by asking you to reach your own verdict on the scale of Section~\ref{sec:01-method} and compare it with Section~\ref{sec:06-verdict}.

\section{Where the evidence stands}
\label{sec:06-verdict}

The published studies point to a clear pattern: \textbf{high-interest currencies have not fallen by their interest advantage on average, so the carry trade has earned a premium, but that premium is paid for with sudden crashes in global downturns.} The table summarizes the evidence as of November 2026.

{\small
\begin{longtable}{@{}>{\raggedright\arraybackslash}p{0.19\linewidth}>{\raggedright\arraybackslash}p{0.13\linewidth}>{\raggedright\arraybackslash}p{\dimexpr 0.45\linewidth-6\tabcolsep\relax}>{\raggedright\arraybackslash}p{0.23\linewidth}@{}}
\toprule
Situation & Holds? & Why & Evidence \\
\midrule
\endfirsthead
\multicolumn{4}{@{}l}{\emph{(continued)}} \\
\toprule
Situation & Holds? & Why & Evidence \\
\midrule
\endhead
\midrule
\multicolumn{4}{r@{}}{\emph{(continued on the next page)}} \\
\endfoot
\bottomrule
\endlastfoot
Fama slope equal to 1 & \textbf{No} & Slopes are below 1 and usually below zero, averaging about $-0.88$ across 75 estimates. & Fama (1984); Froot and Thaler (1990) \\*
\addlinespace[3pt]
Carry portfolios earn nothing & \textbf{No} & Long high-rate, short low-rate portfolios earned substantial returns, largely explained by a carry factor. & Lustig, Roussanov and Verdelhan (2011) \\
\addlinespace[3pt]
Carry in global downturns & \textbf{Yes} & Returns are negatively skewed and lose heavily when volatility rises and funding dries up. & Brunnermeier, Nagel and Pedersen (2008); Menkhoff, Sarno, Schmeling and Schrimpf (2012) \\
\addlinespace[3pt]
Rare disasters not yet seen & \textbf{No} & Carry hedged with currency options still earned a substantial return. & Burnside, Eichenbaum, Kleshchelski and Rebelo (2011) \\
\addlinespace[3pt]
Long horizons and very high inflation & \textbf{Partly} & Slopes are closer to 1 over several years and where large interest gaps match large depreciations. & Later studies; not tested in the lab \\*
\addlinespace[3pt]
Carry after 2008 & \textbf{Unclear} & Narrow interest gaps shrank the carry available; whether the premium per unit of gap changed is open. & Section~\ref{sec:06-stress} \\
\end{longtable}
\addtocounter{table}{-1}% uncaptioned longtable: keep table numbers unchanged
}

The premium looks like payment for risk, not an arbitrage: carry investors provide money and insurance to the rest of the world at the moments when both are scarce.

\textbf{What would change the picture.} The claim would gain support if carry returns since 2008 had been consistently close to zero, as small interest gaps and crowded trades\index{crowding} might imply.

\label{sec:06-open}
\textbf{Open questions.} If the premium is a slowly corrected mispricing\index{mispricing}, it should shrink as more money chases it; has it? And do momentum\index{momentum} or volatility rules sidestep the worst months out of sample, or simply give up the premium?

\section{What practitioners and investors can use}
\label{sec:06-investors}

A currency carry fund earns modest, steady interest differentials in calm years, with most losses concentrated in a few months of global stress, when high-yielding currencies fall together and the borrowed currencies, above all the yen and the Swiss franc, rise. A volatility spike or a margin call\index{margin} forces leveraged traders to sell, which pushes high-yielding currencies down further, and spreading the trade across many pairs helps little because the losses share a common source. Carry tends to lose when stocks fall, so measure its correlation with stocks in bad months, and ask a manager how much leverage\index{leverage} the fund uses and how it did in late 2008 and August 2024.

\begin{callout}
\textbf{What Failed: the yen carry trade, 2008.} Before 2008, borrowing yen at near-zero rates to invest in Australian, New Zealand and other high-yielding currencies was one of the world's most popular trades. As the financial crisis deepened, carry traders rushed to repay their yen loans together, and between July and late October 2008 the Australian and New Zealand dollars lost more than a third of their value against the yen, wiping out years of interest income in weeks. It happened again on a smaller scale in August 2024: after the Bank of Japan raised its policy rate to 0.25\%, the yen rose sharply, carry positions were cut, and on August 5 the Nikkei 225 fell 12.4\%.
\end{callout}
